%%=============================================================================
%% Voorwoord
%%=============================================================================

\chapter*{\IfLanguageName{dutch}{Woord vooraf}{Preface}}%
\label{ch:voorwoord}

Voor u ligt de bachelorproef "Analyse van de sim-to-real uitdagingen bij de integratie van NVIDIA Isaac Sim met ROS 2". Dit werk vormt de afsluiting van mijn opleiding Toegepaste Informatica aan de Hogeschool Gent en is het resultaat van een intensief jaar vol technisch onderzoek, experimenten en diepgaande probleemanalyses op het snijvlak van IT en OT.

Mijn fascinatie voor robotica komt voort uit de unieke synergie tussen code en de fysieke wereld. het moment waarop abstracte algoritmen tastbare gevolgen hebben in de realiteit. Deze passie was de drijfveer om een eigen autonome magazijnrobot (AMR) van de grond af op te bouwen. Het bleek echter een traject met aanzienlijke technische drempels. De transitie van een initiële opzet met een NVIDIA Jetson Nano naar een hybride systeem met een Raspberry Pi 5 (noodzakelijk door de \textit{End-of-Life} status van de Jetson en bijbehorende software compatibiliteitsproblemen) was een van de grootste uitdagingen. Vooral het integreren van de ZED Mini stereoscopische camera voor accurate \textit{Visual Inertial Odometry} (VIO) binnen een gedistribueerd netwerk was een complex vraagstuk waarvoor geen kant-en-klare oplossingen bestonden.

Dit resultaat was niet mogelijk geweest zonder de steun en expertise van verschillende personen en organisaties, die ik hierbij hartelijk wil bedanken.

In de eerste plaats gaat mijn dank uit naar mijn promotor, dhr. S. Verschraege, voor zijn academische begeleiding, het kritische oog op het onderzoeksvlak en zijn bereikbaarheid voor al mijn vragen.

Mijn oprechte dank gaat ook uit naar mijn co-promotor, dhr. W. Martens, en het team bij MR Watts. De open communicatie, het fungeren als klankbord en de technische ondersteuning waren onmisbaar om de complexe versie-mismatches tussen de ZED SDK en ROS 2-wrappers te overwinnen. Tevens dank ik hen voor het ter beschikking stellen van cruciale hardware, zoals de Jetson Nano en de ZED Mini camera.

Daarnaast wil ik de bedrijven Sport-IT en 45Degrees bedanken voor hun financiële steun bij de realisatie van dit project. Een bijzonder woord van dank aan mijn peter en aan Multishop Desomer Plancke voor de vrije toegang tot hun warenhuis, wat een realistische testomgeving bood die essentieel was voor de validatie van de robot.

Op persoonlijk vlak danken ik mijn familie voor hun onuitputtelijke geduld. Zij fungeerden vaak als ``rubber duck'' tijdens technische blokkages en hielpen de leesbaarheid van dit document te verbeteren door hun eerlijke feedback.

Tot slot wil ik de lezer wijzen op de bijbehorende GitHub-repository van deze bachelorproef. Hier zijn alle uitgebreide codefragmenten, configuratiebestanden en bronverwijzingen terug te vinden die in de paper niet volledig aan bod konden komen. Ik hoop dat deze repository en de conclusies in dit werk een solide fundament bieden voor anderen die de uitdagende wereld van \textit{sim-to-real} integratie willen verkennen.

\vspace{1cm}

\noindent Wyffels Robin \
Herne, mei 2026